\documentclass[11pt,a4paper]{article}

% A self-contained manuscript: the bibliography is embedded below.
% Compile: latexmk -pdf speed_monotonicity.tex
\usepackage[T1]{fontenc}
\usepackage[utf8]{inputenc}
\usepackage{newpxtext}
\usepackage{amsmath,amsthm}
\usepackage{newpxmath}
\usepackage[margin=28mm,headheight=14pt]{geometry}
\usepackage{microtype}
\usepackage{mathtools}
\usepackage{aliascnt}
\usepackage{enumitem}
\usepackage{booktabs}
\usepackage{fancyhdr}
\usepackage{titlesec}
\usepackage{hyperref}
\usepackage[nameinlink,noabbrev]{cleveref}
\hypersetup{hidelinks,pdftitle={Speed monotonicity on Galton--Watson trees with at least two offspring},pdfauthor={},pdfsubject={Research draft: tilted moments, local stochastic order, and root response}}
\urlstyle{same}
\setlength{\parindent}{1.1em}
\setlength{\parskip}{2pt}
\setlength{\emergencystretch}{1em}
\linespread{1.04}
\numberwithin{equation}{section}
\titleformat{\section}{\large\bfseries}{\thesection.}{0.6em}{}
\titleformat{\subsection}{\normalsize\bfseries}{\thesubsection.}{0.6em}{}
\titlespacing*{\section}{0pt}{2.0ex plus 0.5ex}{1.0ex}
\titlespacing*{\subsection}{0pt}{1.5ex plus 0.5ex}{0.6ex}
\pagestyle{fancy}
\fancyhf{}
\fancyhead[L]{\small\scshape Speed monotonicity on Galton--Watson trees}
\fancyhead[R]{\small Research draft}
\fancyfoot[C]{\thepage}
\renewcommand{\headrulewidth}{0.25pt}

\theoremstyle{plain}
\newtheorem{theorem}{Theorem}[section]
\newaliascnt{proposition}{theorem}
\newtheorem{proposition}[proposition]{Proposition}
\aliascntresetthe{proposition}
\newaliascnt{lemma}{theorem}
\newtheorem{lemma}[lemma]{Lemma}
\aliascntresetthe{lemma}
\newaliascnt{corollary}{theorem}
\newtheorem{corollary}[corollary]{Corollary}
\aliascntresetthe{corollary}
\theoremstyle{definition}
\newaliascnt{definition}{theorem}
\newtheorem{definition}[definition]{Definition}
\aliascntresetthe{definition}
\theoremstyle{remark}
\newaliascnt{remark}{theorem}
\newtheorem{remark}[remark]{Remark}
\aliascntresetthe{remark}
\newtheorem*{scope}{Scope of the manuscript}
\newtheorem*{lemmaA}{Lemma A (local monotonicity in the escape law)}
\newtheorem*{lemmaB}{Lemma B (strict negativity of the root response)}
\crefname{lemmaA}{Lemma A}{Lemmas A}
\crefname{lemmaB}{Lemma B}{Lemmas B}

\newcommand{\E}{\mathbb E}
\newcommand{\Prob}{\mathbb P}
\newcommand{\R}{\mathbb R}
\newcommand{\N}{\mathbb N}
\newcommand{\Law}{\operatorname{Law}}
\newcommand{\Cov}{\operatorname{Cov}}
\newcommand{\supp}{\operatorname{supp}}
\newcommand{\st}{\mathrm{st}}
\newcommand{\id}{\mathrm{id}}
\newcommand{\V}{\mathcal V}
\newcommand{\J}{\mathcal J}
\newcommand{\one}{\mathbf 1}
\newcommand{\Dini}{D^{+}}
\DeclareMathOperator{\Lip}{Lip}

\title{\vspace{-1.4em}\textbf{Speed monotonicity on Galton--Watson trees\\with at least two offspring}}
\author{}
\date{25 September 2026}

\begin{document}
\maketitle
\vspace{-1.9em}
\begin{center}
\small\textsc{Working research draft}\par
\footnotesize Proof assembled and algebraically checked; independent mathematical review is pending.
\end{center}
\vspace{0.3em}

\begin{abstract}
We give an argument for strict monotonicity of the speed of a biased random
walk on a Galton--Watson tree whose offspring distribution has finite mean
and is supported on $\{2,3,\ldots\}$.  The parameter $\lambda$ is the weight
assigned to the parent edge, so the asserted monotonicity is decreasing on
the transient interval $[0,\E\nu)$.  A two-child recursive distributional
equation yields a strict inequality between three moments under an
inverse-square change of measure.  This scalar inequality implies both
local stochastic-order monotonicity of a symmetrized speed functional
(Lemma A) and strict negativity of its root-only parameter response
(Lemma B).  A finite-difference comparison then proves a strictly negative
upper right Dini derivative for every $1<\lambda<\E\nu$, without
differentiating the infinite-tree escape probability.  A published
small-parameter result completes the interval.  Under an additional
exponential-moment assumption, a published differentiability theorem turns
the Dini-derivative statement into $v'(\lambda)<0$ throughout the transient
regime.  The case with a positive probability of one child is not covered.
\end{abstract}

\noindent\textbf{Keywords.} Galton--Watson tree; biased random walk; speed;
escape probability; recursive distributional equation; stochastic order.

\medskip
\noindent\textbf{Mathematics Subject Classification.} 60J80, 60K37, 60J10.

\section{Introduction and main result}\label{sec:intro}

Let $T$ be a Galton--Watson tree with offspring variable $\nu$.  At a vertex
with $k$ children, a $\lambda$-biased random walk moves to its parent with
probability $\lambda/(\lambda+k)$ and to each child with probability
$1/(\lambda+k)$.  Thus increasing $\lambda$ increases the bias towards the
root.  Throughout the main theorem we assume
\begin{equation}\label{eq:assumptions}
 \Prob(\nu\ge2)=1,
 \qquad m:=\E\nu<\infty.
\end{equation}
The tree is then infinite almost surely.  For $0<\lambda<m$, the distance
from the root divided by time converges almost surely to a deterministic
positive speed $v(\lambda)$; see Lyons, Pemantle and Peres~\cite{LPP96}.
We set $v(0)=1$.

The speed-monotonicity question is whether $v$ decreases when the
parent-edge weight increases.  The pointwise radial drift decreases at
every vertex, but the environment sampled by the walk also changes with
$\lambda$.  Consequently, monotonicity does not follow by averaging local
drifts with respect to the original offspring law.

A\"id\'ekon~\cite{A14} expressed the speed using independent escape
probabilities.  Song, Wang and Xiang~\cite{SWX25} proved strict decrease on
\begin{equation}\label{eq:knowninterval}
 \left[0,\frac{m_1}{1+\sqrt{1-1/m_1}}\right],
 \qquad m_1:=\min\{k:\Prob(\nu=k)>0\}\ge2.
\end{equation}
The upper endpoint of this interval is greater than $1$.  For the
particular law uniform on $\{2,3\}$, the recent preprint of Mandarapu and
Kunkunuru~\cite{MK26} reports a computer-assisted extension to $1.755$.
These results motivate a distributional argument that does not lose its
control when $\lambda$ reaches the minimum offspring number.

\begin{theorem}[Speed monotonicity]\label{thm:main}
Assume~\eqref{eq:assumptions}.  For every $1<\lambda<m$,
\begin{equation}\label{eq:main-dini}
 \Dini v(\lambda)
 :=\limsup_{h\downarrow0}\frac{v(\lambda+h)-v(\lambda)}{h}<0.
\end{equation}
Consequently, $v$ is strictly decreasing on $[0,m)$.
If, in addition, $\E e^{\theta\nu}<\infty$ for some $\theta>0$, then $v$ is
classically differentiable on $(0,m)$ and
\begin{equation}\label{eq:main-derivative}
 v'(\lambda)<0,\qquad 0<\lambda<m.
\end{equation}
\end{theorem}

The extra moment assumption in the last sentence is used only to invoke
the differentiability theorem of Bowditch and Tokushige~\cite{BT20}.
Strict monotonicity and~\eqref{eq:main-dini} use only finite mean.

\begin{scope}
The bias parameter varies while the offspring distribution is fixed.
The theorem is not a claim about monotonicity under stochastic ordering
of two offspring distributions.  No conclusion for $\Prob(\nu=1)>0$ is
asserted.  Lemma A below is local at the relevant escape law, not a global
monotonicity statement on all probability laws on $[0,1]$.
\end{scope}

\subsection{The two comparison lemmas}

For $\lambda>1$ and a probability law $\rho$ on $[0,1]$ other than
$\delta_0$, set
\begin{equation}\label{eq:V-intro}
 \V_\lambda(\rho)
 =\frac{\E_{\rho\otimes\rho}
     [XY/(\lambda-1+X+Y-XY)]}
 {\E_{\rho\otimes\rho}
     [(X+Y-XY)/(\lambda-1+X+Y-XY)]}.
\end{equation}
Let $\pi_\lambda$ denote the law of the quenched probability of escaping
from a vertex without hitting its parent.  We shall derive
$v(\lambda)=\V_\lambda(\pi_\lambda)$ from A\"id\'ekon's formula.

\begin{lemmaA}
Under~\eqref{eq:assumptions}, fix $1<\lambda<m$.  There exist $\eta>0$ and
$\varepsilon>0$ such that, whenever
$|\alpha-\lambda|<\eta$ and $W_1(\rho_i,\pi_\lambda)<\varepsilon$ for
$i=0,1$,
\begin{equation}\label{eq:A-intro}
 \rho_0\le_{\st}\rho_1
 \quad\Longrightarrow\quad
 \V_\alpha(\rho_0)\le\V_\alpha(\rho_1).
\end{equation}
The inequality is strict if $\rho_0\ne\rho_1$.
\end{lemmaA}

\begin{lemmaB}
Under the same assumptions, define
\begin{equation}\label{eq:root-map-intro}
 T_{\lambda,h}(x)
 =\frac{\lambda x}{\lambda+h(1-x)},\qquad
 G_\lambda(h)=\V_{\lambda+h}
       \bigl((T_{\lambda,h})_\#\pi_\lambda\bigr).
\end{equation}
Then $G_\lambda$ is differentiable at $0$ and $G_\lambda'(0)<0$.
\end{lemmaB}

Lemma A compares nearby laws at a fixed parameter.  Lemma B changes the
parameter and only the root contribution to the escape probability.  The
full RDE response is smaller than this root-only comparison, in stochastic
order.  Joint local uniformity in $\alpha$ and $\rho$ is essential when
combining the two lemmas.

The main new input is a scalar tilted-moment inequality
$\Gamma_\lambda(\pi_\lambda)>0$.  We prove it in
Section~\ref{sec:moments}; Section~\ref{sec:functional} proves that this
single inequality implies both comparison lemmas.  Section~\ref{sec:closure}
then supplies the finite-difference and continuity arguments.  In
particular, no formal differentiation of an infinite-tree fixed point is
used.

\section{Escape probabilities and the symmetrized speed}\label{sec:setup}

\subsection{The canonical fixed point}

Attach an artificial parent $e_*$ to the root $e$.  For a fixed tree, let
\begin{equation}\label{eq:beta}
 \beta_\lambda(T)=P_e^T(\tau_{e_*}=\infty),
 \qquad \pi_\lambda=\Law(\beta_\lambda(T)).
\end{equation}
We always use this canonical escape probability, rather than an arbitrary
solution of a recursive equation.  In the regime $\lambda<m$, every
subtree survives and is transient almost surely, and hence
$0<\beta_\lambda<1$ almost surely.  The strict upper bound also follows
from the positive chance of making the first step to the parent.

If $\beta_{\lambda,i}$ are the escape probabilities in the child subtrees,
a first-step calculation gives
\begin{equation}\label{eq:rde}
 \beta_\lambda
 =F_\lambda(S_\lambda),\qquad
 F_\lambda(s)=\frac{s}{\lambda+s},\qquad
 S_\lambda=\sum_{i=1}^{\nu}\beta_{\lambda,i}.
\end{equation}
Under the Galton--Watson law the child variables are independent, have law
$\pi_\lambda$, and are independent of $\nu$.

For clarity, the canonical solution can also be obtained by assigning
boundary value $1$ at generation $n$ and iterating $F_\lambda$ towards the
root.  These finite-depth values decrease as $n$ increases and converge
to~\eqref{eq:beta}.  Since $F_\lambda$ is increasing in each child input
and decreasing in $\lambda$, induction and passage to the limit give
\begin{equation}\label{eq:pathwise-monotone}
 \beta_{\lambda+h}(T)\le\beta_\lambda(T),\qquad h>0.
\end{equation}
All comparisons between different parameters use this same-tree coupling.

\subsection{An algebraic symmetrization}

\begin{proposition}[Symmetric speed formula]\label{prop:symmetric}
Under~\eqref{eq:assumptions}, for $1<\lambda<m$,
\begin{equation}\label{eq:symmetric-speed}
 v(\lambda)=\V_\lambda(\pi_\lambda).
\end{equation}
\end{proposition}
\begin{proof}
Write $c=\lambda-1$ and let $X_0,X_1,\ldots$ be independent with law
$\pi_\lambda$, independent of $\nu$.  Set $S=\sum_{i=1}^{\nu}X_i$.
A\"id\'ekon's formula~\cite[Theorem 1.1]{A14} is
\begin{equation}\label{eq:aidekon}
 v(\lambda)
 =\frac{\E[(\nu-\lambda)X_0/(c+X_0+S)]}
        {\E[(\nu+\lambda)X_0/(c+X_0+S)]}.
\end{equation}
Conditional on $\nu=k$, the $k+1$ inputs are exchangeable.  Therefore
\begin{equation}\label{eq:exchange}
 \E\frac{\nu X_0}{c+X_0+S}
 =\E\frac S{c+X_0+S}.
\end{equation}
Let $Y=S/(\lambda+S)$.  By the RDE, $Y$ has law $\pi_\lambda$ and is
independent of $X_0$.  Since
\begin{equation}\label{eq:algebra-S}
 S=\frac{\lambda Y}{1-Y},\qquad
 c+X_0+S=\frac{c+X_0+Y-X_0Y}{1-Y},
\end{equation}
the numerator of~\eqref{eq:aidekon} becomes
\begin{equation}
 \lambda\E\frac{Y-X_0+X_0Y}{c+X_0+Y-X_0Y}
 =\lambda\E\frac{X_0Y}{c+X_0+Y-X_0Y}.
\end{equation}
The antisymmetric term has expectation zero.  The denominator becomes
$\lambda\E[(X_0+Y-X_0Y)/(c+X_0+Y-X_0Y)]$.  Cancelling $\lambda$ proves
the result.
\end{proof}

\subsection{Continuity and elementary reweighting}

We use the following published continuity input, separately from any
speed differentiability result.

\begin{proposition}[Continuity input]\label{prop:continuity}
Under~\eqref{eq:assumptions}, the map
$\lambda\mapsto\E\beta_\lambda$ is continuous on $(0,m)$.
Consequently, $\lambda\mapsto\pi_\lambda$ is continuous in $W_1$, and
$v$ is continuous on $(1,m)$.
\end{proposition}
\begin{proof}
Continuity of the annealed escape probability is
\cite[Proposition 2.9]{BT20}.  Its proof in the no-leaf case uses
\cite[Lemma 4.7]{BT20}; these continuity statements do not require the
exponential-moment hypothesis used for differentiability in
\cite[Theorem 1.1]{BT20}.

For parameters $\alpha,\lambda$, the same-tree coupling in
\eqref{eq:pathwise-monotone} is ordered.  Hence
\begin{equation}\label{eq:w1-cont}
 W_1(\pi_\alpha,\pi_\lambda)
 =\bigl|\E\beta_\alpha-\E\beta_\lambda\bigr|.
\end{equation}
Indeed, this coupling gives the upper bound, and the $1$-Lipschitz test
function $x\mapsto x$ gives the reverse inequality.  For parameters in a
compact subinterval of $(1,m)$, the two kernels in
\eqref{eq:V-intro} are bounded and jointly continuous, and the denominator
is positive.  Formula~\eqref{eq:symmetric-speed} proves continuity of $v$.
\end{proof}

\begin{lemma}[Increasing reweighting]\label{lem:tilt}
Let $M$ be a probability measure on an interval.  Suppose $r\ge0$ is
increasing and $0<\int r\,dM<\infty$.  For every bounded increasing
function $f$,
\begin{equation}\label{eq:tilt}
 \frac{\int fr\,dM}{\int r\,dM}\ge\int f\,dM.
\end{equation}
For a decreasing $r$, the inequality is reversed.
\end{lemma}
\begin{proof}
For independent $U,V$ with law $M$,
\begin{equation}
 2\Cov_M(f,r)
 =\E[(f(U)-f(V))(r(U)-r(V))].
\end{equation}
The integrand is nonnegative if both functions are increasing and
nonpositive if their monotonicities are opposite.  Dividing by
$\int r\,dM$ proves the claim.  The same argument applies by truncation
whenever the displayed expectations exist.
\end{proof}

\section{The two-child tilted-moment inequality}\label{sec:moments}

This section is an independent statement about RDE fixed points.  No
random-walk regularity theorem is used here.

Fix $\lambda>1$, write $c=\lambda-1$, and let $X\in(0,1)$ satisfy
\begin{equation}\label{eq:abstract-rde}
 X\overset d=\frac{S}{\lambda+S},\qquad
 S=\sum_{i=1}^{\nu}X_i,\qquad \nu\ge2,
\end{equation}
where $\nu$ is finite almost surely, the $X_i$ are independent copies of
$X$, and they are independent of $\nu$.  Define
\begin{equation}\label{eq:tilted-moments}
 a=\E(c+X)^{-2},\qquad
 dM=\frac{(c+X)^{-2}}a\,d\Prob,
 \qquad
 \mu=\E X,\quad k=E_MX,\quad j=E_MX^2,\quad \ell=E_MX^3.
\end{equation}
All these quantities are finite, and $0<\mu,k<1$.

\begin{lemma}[Two-child constraint]\label{lem:two-child}
Under~\eqref{eq:abstract-rde},
\begin{equation}\label{eq:moment-constraint}
 \lambda(k-j)\ge2k(1-2k+j).
\end{equation}
Equivalently,
\begin{equation}\label{eq:j-upper}
 k^2\le j\le\bar j:=\frac{k(c-1+4k)}{c+1+2k}.
\end{equation}
In particular,
\begin{equation}\label{eq:d-positive}
 d:=c+2k-1\ge0.
\end{equation}
\end{lemma}
\begin{proof}
On the root environment $(\nu,X_1,X_2,\ldots)$ define
\begin{equation}\label{eq:Q-root}
 dQ=\frac{(c+S)^{-2}}{\E(c+S)^{-2}}\,d\Prob.
\end{equation}
Write $S=X_1+H$, where $H=\sum_{i=2}^{\nu}X_i$.  The variable $X_1$ is
independent of $H$.  The $X_1$-marginal under $Q$, relative to $M$, has
density proportional to
\begin{equation}\label{eq:R-tilt}
 R(x)=(c+x)^2\E\frac1{(c+x+H)^2}
     =\E\left(\frac{c+x}{c+x+H}\right)^2.
\end{equation}
This is increasing in $x$.  By~\cref{lem:tilt}, $E_QX_1\ge k$.
Conditional exchangeability of the children gives
\begin{equation}\label{eq:two-children}
 E_QS=E_Q[\nu X_1]\ge2E_QX_1\ge2k.
\end{equation}
The first equality follows by conditioning on every value of $\nu$;
the tilt is symmetric in all the child inputs.  These quantities are
integrable even without $\E\nu<\infty$, since
$S/(c+S)^2\le1/(4c)$.

The RDE implies
\begin{equation}
 S\overset d=\frac{\lambda X}{1-X},\qquad
 c+S\overset d=\frac{c+X}{1-X}.
\end{equation}
Consequently,
\begin{equation}\label{eq:EQ-ratio}
 E_QS
 =\lambda\frac{E_M[X(1-X)]}{E_M[(1-X)^2]}
 =\lambda\frac{k-j}{1-2k+j}.
\end{equation}
The denominator is strictly positive.  Combining this identity with
\eqref{eq:two-children} yields~\eqref{eq:moment-constraint} and its upper
bound for $j$.  The lower bound $j\ge k^2$ is Jensen's inequality.
Finally,
\begin{equation}\label{eq:jgap}
 \bar j-k^2
 =\frac{k(1-k)(c+2k-1)}{c+1+2k}.
\end{equation}
Its nonnegativity gives~\eqref{eq:d-positive}.
\end{proof}

\begin{proposition}[Strict scalar inequality]\label{prop:Gamma}
Under~\eqref{eq:abstract-rde},
\begin{equation}\label{eq:Gamma}
 \Gamma:=\lambda(2k-\mu)-(1-\mu)(k-j)>0.
\end{equation}
\end{proposition}
\begin{proof}
If $k\ge1/2$, rewrite
\begin{equation}\label{eq:Gamma-largek}
 \Gamma=\lambda(2k-1)+(1-\mu)(\lambda-k+j)>0.
\end{equation}
The second term is strictly positive because $\mu<1$ and
$\lambda-k+j>0$.

Suppose henceforth that $0<k<1/2$.  Reversing the change of measure in
\eqref{eq:tilted-moments} gives
\begin{equation}\label{eq:unweight}
 \mu=\frac{c^2k+2cj+\ell}{c^2+2ck+j}
 \le\frac{c^2k+(2c+1)j}{c^2+2ck+j}=:\mu_+(j),
\end{equation}
since $X^3\le X^2$ under $M$.  The expression in~\eqref{eq:Gamma} is
strictly decreasing in $\mu$, its coefficient being
$-(\lambda-k+j)<0$.  It therefore suffices to prove
\begin{equation}\label{eq:F-def}
 F(j):=(c^2+2ck+j)
 \bigl[(c+1)(2k-\mu_+(j))-(1-\mu_+(j))(k-j)\bigr]>0.
\end{equation}
This is a quadratic polynomial in $j$, with leading coefficient $-2c$.
It is concave.  On the interval in~\eqref{eq:j-upper}, it is therefore
enough to check its two endpoints.

Put $d=c+2k-1\ge0$.  Direct polynomial expansion yields
\begin{equation}\label{eq:F-endpoints}
 F(k^2)=kP_3(d,k),\qquad
 F(\bar j)=\frac{kP_5(d,k)}{(d+2)^2},
\end{equation}
where
\begin{align}\label{eq:P3}
 P_3(d,k)
 ={}&d^3+(1-k)(k+3)d^2\notag\\
 &+(1-k)(3-2k-2k^2)d+(1-k)^2(1-2k),
\end{align}
and
\begin{align}\label{eq:P5}
 P_5(d,k)
 ={}&d^5+2(3-k)d^4+2(6-k-3k^2)d^3\notag\\
 &+2(1-k)(5+12k-8k^2)d^2\notag\\
 &+2(1-k)(3+7k-16k^2+4k^3)d\notag\\
 &+4(1-k)^2(1-2k).
\end{align}
Every coefficient displayed in~\eqref{eq:P3}--\eqref{eq:P5} is positive
for $0<k<1/2$.  For completeness,
$3-2k-2k^2>3/2$,
$6-k-3k^2>19/4$, and $5+12k-8k^2>0$.  The remaining cubic coefficient
has the useful factorization
\begin{equation}\label{eq:cubic-coeff}
 3+7k-16k^2+4k^3=3+k(1-2k)(7-2k)>0.
\end{equation}
The constant coefficients in both polynomials are strictly positive.
Thus $P_3(d,k)>0$ and $P_5(d,k)>0$ whenever $d\ge0$ and $0<k<1/2$.
Concavity now gives $F(j)>0$ throughout the allowed interval.  By
\eqref{eq:unweight}--\eqref{eq:F-def}, this proves $\Gamma>0$.
\end{proof}

\begin{remark}\label{rem:exact-algebra}
The proof uses only the two explicit identities~\eqref{eq:F-endpoints}
and elementary positivity of their coefficients.  The accompanying
\texttt{verification/check\_algebra.py} checks the identities exactly
using rational symbolic arithmetic.  No numerical quadrature, sampling,
or parameter grid is part of the proof.
\end{remark}

\section{One scalar condition implies Lemmas A and B}\label{sec:functional}

We now temporarily forget the RDE.  Fix $\lambda>1$ and any probability
law $\pi$ supported on $(0,1)$.  All unmarked expectations in this section
use independent $X,Y\sim\pi$.  Set
\begin{align}\label{eq:general-functional}
 c&=\lambda-1,&
 t(x,y)&=x+y-xy,&
 A(x,y)&=c+t(x,y),\notag\\
 N&=\E\frac{XY}{A(X,Y)},&
 Z&=\E\frac{t(X,Y)}{A(X,Y)},&
 v&=\frac NZ.
\end{align}
Then $Z>0$ and $0<v<1$.  Define $\mu,k,j$ by
\eqref{eq:tilted-moments}, and assume for now that $\Gamma>0$ as defined
in~\eqref{eq:Gamma}.

For a test function $f$ let
\begin{equation}\label{eq:J-operator}
 (\J f)(x)=\E\frac{f(Y)}{A(x,Y)^2},\qquad 0\le x\le1.
\end{equation}
Also set
\begin{align}\label{eq:h-f-kappa}
 h_v(y)&=(1+v)y-v,\notag\\
 f_v(y)&=h_v(y)-\frac{1-v}{2\lambda}y(1-y),\notag\\
 \kappa_v(y)&=y^2+c h_v(y),\qquad
 \Psi=\J\kappa_v.
\end{align}
The quantities $v,h_v,f_v$ in derivative identities below are evaluated
at the base point.  When differentiating $N-vZ$, the scalar $v$ is held
fixed.

\subsection{A positive integral kernel}

\begin{lemma}[Kernel positivity from $\Gamma$]\label{lem:kernel}
If $\Gamma>0$, then, for all $x\in[0,1]$,
\begin{align}\label{eq:Jf-positive}
 (\J f_v)(x)&\ge\frac{\Gamma}{\lambda^3(2-\mu)}>0,\\
 \Psi(x)&\ge\frac{c\Gamma}{\lambda^3(2-\mu)}>0.
 \label{eq:Psi-positive}
\end{align}
\end{lemma}
\begin{proof}
Let $W=\E[X/A(X,Y)]$.  Symmetry gives $Z+N=2W$.  Conditional on $X=x$,
the function $y\mapsto A(x,y)^{-1}$ is decreasing.  By
\cref{lem:tilt},
\begin{equation}
 \E\frac{Y}{A(x,Y)}\le\mu\E\frac1{A(x,Y)}.
\end{equation}
Multiplying by $x$ and integrating gives $N\le\mu W$, and hence
\begin{equation}\label{eq:v-mu}
 v\le\frac{\mu}{2-\mu}.
\end{equation}
For $0\le y\le1$,
\begin{align}\label{eq:f-monotone}
 \partial_v f_v(y)&=-(1-y)\left(1-\frac{y}{2\lambda}\right)\le0,\notag\\
 f_v'(y)&=1+v-\frac{1-v}{2\lambda}(1-2y)>0.
\end{align}
Therefore
\begin{align}\label{eq:f-tilted-average}
 E_M f_v(X)
 &\ge E_M f_{\mu/(2-\mu)}(X)\notag\\
 &=\frac{\lambda(2k-\mu)-(1-\mu)(k-j)}{\lambda(2-\mu)}
 =\frac{\Gamma}{\lambda(2-\mu)}.
\end{align}

For fixed $x\in[0,1]$, tilt the law of $Y$ by $A(x,Y)^{-2}$.  Relative
to $M$, its density is proportional to
\begin{equation}
 r_x(y)=\left(\frac{c+y}{c+x+(1-x)y}\right)^2.
\end{equation}
This function is increasing, because
\begin{equation}\label{eq:likelihood-deriv}
 \frac{d}{dy}\frac{c+y}{c+x+(1-x)y}
 =\frac{\lambda x}{A(x,y)^2}\ge0.
\end{equation}
Apply~\cref{lem:tilt} to the increasing function $f_v$.  Since
$A(x,y)\le\lambda$ on $[0,1]^2$,~\eqref{eq:f-tilted-average} yields
\eqref{eq:Jf-positive}.  Finally,
\begin{equation}
 \kappa_v(y)=y^2+c f_v(y)+\frac{c(1-v)}{2\lambda}y(1-y),
\end{equation}
and the two terms besides $c f_v$ are nonnegative.  This proves
\eqref{eq:Psi-positive}.
\end{proof}

\subsection{Lemma A: a genuinely local statement}

\begin{proposition}[Local stochastic-order comparison]\label{prop:A}
Fix $\lambda>1$ and a law $\pi$ supported on $(0,1)$ with
$\Gamma_\lambda(\pi)>0$.  Then the conclusion of Lemma A holds, with
$\pi$ in place of $\pi_\lambda$.
\end{proposition}
\begin{proof}
For any parameter $\alpha>1$ and law $\rho\ne\delta_0$, define
$v_{\alpha,\rho}=\V_\alpha(\rho)$ and
\begin{equation}\label{eq:general-Psi}
 \Psi_{\alpha,\rho}(x)
 =\int\frac{y^2+(\alpha-1)((1+v_{\alpha,\rho})y-v_{\alpha,\rho})}
              {(\alpha-1+x+y-xy)^2}\,\rho(dy).
\end{equation}
The map $(\alpha,\rho)\mapsto\Psi_{\alpha,\rho}$ is continuous at
$(\lambda,\pi)$ in the uniform norm on $[0,1]$, when laws are given the
$W_1$ topology.  To see this, restrict $\alpha$ to a compact subinterval
of $(1,\infty)$.  All relevant kernels are jointly continuous on the
corresponding compact set, with denominators bounded away from zero.
Product integrals depend continuously on $\rho$, and $Z_\lambda(\pi)>0$
keeps the ratio defining $v_{\alpha,\rho}$ well defined nearby.  Uniform
continuity in $x$ gives the asserted uniform-norm continuity.

By~\eqref{eq:Psi-positive}, there exist $\eta,\varepsilon,\delta>0$ such
that
\begin{equation}\label{eq:uniform-neighborhood}
 \Psi_{\alpha,\rho}(x)\ge\delta
 \quad\text{for every }x\in[0,1]
\end{equation}
whenever $|\alpha-\lambda|<\eta$ and $W_1(\rho,\pi)<\varepsilon$.
Shrink $\varepsilon$ if necessary so these laws do not include $\delta_0$.

Suppose $\rho_0\le_{\st}\rho_1$ are in this ball.  Let $U$ be uniform
on $(0,1)$, and couple $X_i=F_{\rho_i}^{-1}(U)$, so $X_0\le X_1$.
Put $X_s=(1-s)X_0+sX_1$ and $\rho_s=\Law(X_s)$.  Quantile interpolation
gives
\begin{equation}\label{eq:quantile-ball}
 W_1(\rho_s,\pi)
 \le(1-s)W_1(\rho_0,\pi)+sW_1(\rho_1,\pi)<\varepsilon.
\end{equation}
Let $Y_s$ be an independent copy of this coupled interpolation.  Direct
differentiation under the bounded kernels gives
\begin{equation}\label{eq:Gateaux}
 \frac{d}{ds}\V_\alpha(\rho_s)
 =\frac{2}{Z_\alpha(\rho_s)}
   \E[(X_1-X_0)\Psi_{\alpha,\rho_s}(X_s)].
\end{equation}
For completeness, the single-coordinate calculation behind this formula
is
\begin{equation}\label{eq:single-variation}
 \frac{\partial}{\partial x}
 \frac{xy-v(x+y-xy)}{\alpha-1+x+y-xy}
 =\frac{y^2+(\alpha-1)((1+v)y-v)}
        {(\alpha-1+x+y-xy)^2},
\end{equation}
with $v$ held fixed.  The quotient rule at each $s$ uses
$v=v_{\alpha,\rho_s}$; exchangeability supplies the factor $2$.

Equations~\eqref{eq:uniform-neighborhood} and~\eqref{eq:Gateaux} show that
the derivative is nonnegative.  If $\rho_0\ne\rho_1$, then
$\E(X_1-X_0)>0$ under this ordered coupling, so the derivative is strictly
positive.  Integrating in $s$ proves the proposition.
\end{proof}

\subsection{Lemma B: the root-response cancellation}

Set
\begin{equation}\label{eq:root-response-quantities}
 b(x)=\frac{x(1-x)}{\lambda},\qquad
 \mathsf C=\E\frac{v\,t(X,Y)-XY}{A(X,Y)^2},\qquad
 \Delta=2\E[b(X)\Psi(X)]-\mathsf C.
\end{equation}

\begin{lemma}[Exact cancellation]\label{lem:cancellation}
For every law $\pi$ supported on $(0,1)$ and every $\lambda>1$,
\begin{equation}\label{eq:Delta-identity}
 \Delta=2\E[X(1-X)(\J f_v)(X)].
\end{equation}
\end{lemma}
\begin{proof}
For $x,y\in[0,1]$, the following polynomial identity holds:
\begin{align}\label{eq:pointwise-cancellation}
 &x(1-x)\kappa_v(y)+y(1-y)\kappa_v(x)
          +\lambda\{xy-vt(x,y)\}\notag\\
 &\qquad=
 \lambda\{x(1-x)f_v(y)+y(1-y)f_v(x)\}
          +A(x,y)\{xy-vt(x,y)\}.
\end{align}
It follows by expanding the definitions in~\eqref{eq:h-f-kappa}.
Divide by $\lambda A(x,y)^2$ and integrate.  The last term on the right
has expectation zero because $N-vZ=0$.  Symmetry of $X,Y$ identifies the
remaining two sides with~\eqref{eq:root-response-quantities} and
\eqref{eq:Delta-identity}.
\end{proof}

\begin{proposition}[Negative root response]\label{prop:B}
Fix $\lambda>1$ and a law $\pi$ supported on $(0,1)$ with
$\Gamma_\lambda(\pi)>0$.  Define
\begin{equation}
 G(h)=\V_{\lambda+h}((T_{\lambda,h})_\#\pi),\qquad
 T_{\lambda,h}(x)=\frac{\lambda x}{\lambda+h(1-x)}.
\end{equation}
Then $G$ is differentiable at $0$ and
\begin{equation}\label{eq:root-negative}
 G'(0)=-\frac\Delta Z
 \le-\frac{2\Gamma\E[X(1-X)]}{Z\lambda^3(2-\mu)}<0.
\end{equation}
\end{proposition}
\begin{proof}
For $h$ in a sufficiently small interval around $0$, $\lambda+h>1$ and
$T_{\lambda,h}$ maps $[0,1]$ into itself.  All parameter and spatial
derivatives involved are bounded on this compact set, so differentiation
under the integrals is justified.  The transport velocity is
\begin{equation}\label{eq:root-velocity}
 \left.\partial_hT_{\lambda,h}(x)\right|_{h=0}=-b(x).
\end{equation}
The explicit parameter derivative of $N-vZ$, holding $v$ and the law
fixed, equals $\mathsf C$.  Combining this with the two transported
coordinates and~\eqref{eq:single-variation} gives
\begin{equation}
 G'(0)=\frac{\mathsf C-2\E[b(X)\Psi(X)]}{Z}=-\frac\Delta Z.
\end{equation}
Use~\cref{lem:cancellation} and~\eqref{eq:Jf-positive} to obtain the
quantitative bound in~\eqref{eq:root-negative}.  It is strict because
$\E[X(1-X)]>0$.
\end{proof}

\begin{corollary}[The two lemmas at the RDE fixed point]\label{cor:AB}
Under~\eqref{eq:assumptions}, both Lemma A and Lemma B hold for every
$1<\lambda<m$.
\end{corollary}
\begin{proof}
Apply~\cref{prop:Gamma} to $X=\beta_\lambda$, and then apply
\cref{prop:A,prop:B}.  The RDE is used only to establish $\Gamma>0$;
the two comparison propositions themselves are distributional statements.
\end{proof}

\section{Finite-difference closure and the speed theorem}\label{sec:closure}

\subsection{A pathwise comparison retaining the parameter change}

\begin{lemma}[Freezing the descendants]\label{lem:root-comparison}
For $h>0$ with $\lambda+h<m$, the same-tree coupling satisfies
\begin{equation}\label{eq:root-comparison}
 \beta_{\lambda+h}\le T_{\lambda,h}(\beta_\lambda).
\end{equation}
Consequently,
\begin{equation}\label{eq:law-root-comparison}
 \pi_{\lambda+h}\le_{\st}(T_{\lambda,h})_\#\pi_\lambda.
\end{equation}
\end{lemma}
\begin{proof}
Every child escape probability decreases with the parameter, so
$S_{\lambda+h}\le S_\lambda$.  Hence
\begin{equation}
 \beta_{\lambda+h}
 =\frac{S_{\lambda+h}}{\lambda+h+S_{\lambda+h}}
 \le\frac{S_\lambda}{\lambda+h+S_\lambda}.
\end{equation}
Solving the base-point RDE for $S_\lambda$ gives
$S_\lambda=\lambda\beta_\lambda/(1-\beta_\lambda)$.  Substitution
identifies the last expression with $T_{\lambda,h}(\beta_\lambda)$.
\end{proof}

\subsection{Negative upper right Dini derivative}

\begin{proof}[Proof of~\eqref{eq:main-dini}]
Fix $1<\lambda<m$ and abbreviate $\pi=\pi_\lambda$.  Take the joint
parameter-law neighborhood supplied by~\cref{prop:A} at $(\lambda,\pi)$.
By~\cref{prop:continuity},
$W_1(\pi_{\lambda+h},\pi)\to0$.  Also
\begin{equation}
 0\le x-T_{\lambda,h}(x)
 =\frac{h x(1-x)}{\lambda+h(1-x)}\le\frac h{4\lambda},
\end{equation}
so $(T_{\lambda,h})_\#\pi$ converges to $\pi$ in $W_1$ as $h\downarrow0$.
Thus both laws in~\eqref{eq:law-root-comparison}, as well as their
quantile interpolation, lie in the same local neighborhood for all
sufficiently small $h>0$.  Lemma A can be applied at the \emph{perturbed}
parameter $\lambda+h$, not merely at $\lambda$.  It yields
\begin{align}\label{eq:closing-comparison}
 v(\lambda+h)
 &=\V_{\lambda+h}(\pi_{\lambda+h})\notag\\
 &\le\V_{\lambda+h}((T_{\lambda,h})_\#\pi)
 =G_\lambda(h).
\end{align}
Since $G_\lambda(0)=v(\lambda)$, Lemma B implies
\begin{equation}\label{eq:quantitative-dini}
 \Dini v(\lambda)
 \le G_\lambda'(0)
 \le-\frac{2\Gamma_\lambda(\pi_\lambda)
                  \E[\beta_\lambda(1-\beta_\lambda)]}
 {Z_\lambda(\pi_\lambda)\lambda^3(2-\E\beta_\lambda)}<0.
\end{equation}
Every differentiation in this argument concerns bounded finite-dimensional
kernels with the base-point law held fixed.  In particular, the proof
has not used $\partial_\lambda\beta_\lambda$, or differentiation of its law.
\end{proof}

\begin{lemma}[A Dini-derivative criterion]\label{lem:dini}
Let $f$ be continuous on an interval $I$.  If
\begin{equation}
 \limsup_{h\downarrow0}\frac{f(x+h)-f(x)}h<0
\end{equation}
for every $x$ in the interior of $I$, then $f$ is strictly decreasing
there.
\end{lemma}
\begin{proof}
Suppose $a<b$ are interior points with $f(b)\ge f(a)$.  Continuity gives
a point where $f$ attains its minimum on $[a,b]$, and such a point can
be chosen strictly before $b$.  Indeed, if $b$ attains the minimum, then
$f(b)\ge f(a)$ forces $a$ to attain it as well.  At a minimizing point
$x<b$, all sufficiently small right difference quotients are nonnegative.
This contradicts the assumed strict negativity.
\end{proof}

\subsection{Completing the interval and the differentiability statement}

\begin{proof}[Completion of~\cref{thm:main}]
Continuity from~\cref{prop:continuity}, together with
\eqref{eq:main-dini} and~\cref{lem:dini}, gives strict decrease on $(1,m)$.
The published result~\eqref{eq:knowninterval} gives strict decrease on a
closed interval whose upper endpoint is strictly greater than $1$;
see~\cite{SWX25}, Theorem 1.2 in its arXiv version.  The two intervals
overlap, so their conclusions combine to give strict decrease on
$[0,m)$.  In particular there is no gap at $\lambda=1$.

If $\E e^{\theta\nu}<\infty$ for some $\theta>0$,
\cite[Theorem 1.1]{BT20} gives classical differentiability throughout
$(0,m)$ in the present no-leaf, no-unary setting.  At $\lambda>1$,
\eqref{eq:quantitative-dini} then implies $v'(\lambda)<0$.  On the
small-parameter interval, the strict derivative inequality is explicitly
proved in~\cite{SWX25}, Lemma 2.5 of the arXiv version.  This supplies
$0<\lambda\le1$ as well and proves~\eqref{eq:main-derivative}.
\end{proof}

\begin{remark}[What finite mean does and does not supply]
Without the additional regularity assumption, the conclusion proved
here is strict monotonicity together with the negative Dini bound
\eqref{eq:quantitative-dini}.  We do not infer existence of a classical
derivative merely from strict monotonicity.  Whenever that derivative
exists at a point $\lambda>1$, it is strictly negative by the same bound.
\end{remark}

\section{Checks, limitations, and proof provenance}\label{sec:scope}

\subsection{The regular-tree check}

If $\nu\equiv d\ge2$, then
\begin{equation}
 \beta_\lambda=1-\frac\lambda d,
 \qquad
 v(\lambda)=\frac{d-\lambda}{d+\lambda}.
\end{equation}
The tilted law is still a point mass, so
$\mu=k=\beta_\lambda$, $j=\beta_\lambda^2$, and
\begin{equation}
 \Gamma=\beta_\lambda
       \bigl(\lambda-(1-\beta_\lambda)^2\bigr)>0
 \qquad(1<\lambda<d).
\end{equation}
For a point mass at $x$, $\V_\lambda(\delta_x)=x/(2-x)$, independently of
the explicit parameter.  Accordingly, the root-only derivative and the
full speed derivative are respectively
\begin{equation}\label{eq:regular-check}
 G_\lambda'(0)=-\frac{2(d-\lambda)}{(d+\lambda)^2},
 \qquad
 v'(\lambda)=-\frac{2d}{(d+\lambda)^2}.
\end{equation}
Both are negative, and the full derivative is the more negative one,
as predicted by~\eqref{eq:closing-comparison}.  Their different orders
near $\lambda=d$ also explain why the proof compares finite differences
rather than identifying root and full derivatives.

\subsection{The precise obstruction when one child is allowed}

The restriction $\nu\ge2$ enters in~\eqref{eq:two-children}.  If instead
$\Prob(\nu\ge1)=1$ and $p=\Prob(\nu=1)>0$, the same argument only gives
\begin{equation}\label{eq:unary-constraint}
 \lambda(k-j)\ge2k(1-2k+j)-pk.
\end{equation}
Indeed,
$\nu\ge2-\one_{\{\nu=1\}}$, and under~\eqref{eq:Q-root},
\begin{equation}
 E_Q[X_1\one_{\{\nu=1\}}]
 =\frac{p\,\E[X/(c+X)^2]}{\E[(c+S)^{-2}]}
 =\frac{pk}{1-2k+j}.
\end{equation}
Thus the upper bound becomes
\begin{equation}\label{eq:unary-jupper}
 j\le\frac{k(c-1+p+4k)}{c+1+2k},
\end{equation}
and feasibility with $j\ge k^2$ gives only
\begin{equation}
 (1-k)(c+2k-1)+p\ge0.
\end{equation}
The crucial condition $d=c+2k-1\ge0$ need not follow.  The coefficient
proof in~\cref{prop:Gamma} therefore does not extend.  In addition,
Lemma A in this manuscript is deduced from $\Gamma>0$ and is asserted
only in the scope stated above.  None of these observations establishes
the sign of a pipe-resummed response in the unary case.

\subsection{Bias monotonicity versus offspring ordering}

The coupling in~\eqref{eq:root-comparison} uses one fixed offspring law and
changes $\lambda$.  It does not compare speeds for two different offspring
laws.  Such a comparison would additionally require an appropriate
continuous path of fixed points for the varying offspring distributions,
and a proof that the required local comparisons can be connected along
that path.  No such theorem is included here.  In particular,
$\nu$ denotes the offspring variable throughout, whereas $v$ denotes the
speed.

\subsection{What has been checked}

The analytic proof consists of the reweighting argument in
\cref{lem:two-child}, the scalar polynomial argument in
\cref{prop:Gamma}, the kernel and cancellation identities in
Section~\ref{sec:functional}, and the local finite-difference closure in
Section~\ref{sec:closure}.  The bibliography supplies the speed formula,
annealed escape continuity, a low-parameter theorem, and the optional
classical differentiability theorem.  These inputs are not asserted as
new results.

The companion script performs exact symbolic checks of
\eqref{eq:F-endpoints}, \eqref{eq:pointwise-cancellation}, the transport
velocity, the moment rearrangements, and the regular-tree check.
The positive-polynomial argument is printed in full, so the proof does
not rely on a numerical certificate.  The script does not formally
verify conditional independence, differentiation under integrals, or the
use of published results; those arguments are given in the text.  This
working draft has not been independently peer reviewed, and no claim of
externally established priority is made.

\appendix
\section{A quantitative distributional criterion}\label{app:criterion}

The main argument proves a statement that may be useful separately from
the tree model.

\begin{theorem}[Distributional criterion]\label{thm:distributional}
Let $\lambda>1$, let $\pi$ be supported on $(0,1)$, and define
$\mu,k,j,\Gamma$ as in~\eqref{eq:tilted-moments} and~\eqref{eq:Gamma}.
If $\Gamma>0$, then
\begin{equation}
 \inf_{x\in[0,1]}\Psi_{\lambda,\pi}(x)
 \ge\frac{(\lambda-1)\Gamma}{\lambda^3(2-\mu)}>0,
\end{equation}
and
\begin{equation}
 \left.\frac d{dh}\V_{\lambda+h}
       ((T_{\lambda,h})_\#\pi)\right|_{h=0}
 \le-\frac{2\Gamma\E_\pi[X(1-X)]}
 {Z_\lambda(\pi)\lambda^3(2-\mu)}<0.
\end{equation}
Moreover, the first conclusion implies local strict stochastic-order
monotonicity jointly in a neighborhood of the parameter and the law.
\end{theorem}
\begin{proof}
This is exactly~\cref{lem:kernel,prop:A,prop:B}.  Those results use no
fixed-point assumption.
\end{proof}

The role of the Galton--Watson structure is consequently isolated:
independence of the first child from the remaining children and the
presence of at least two children imply $\Gamma>0$.  Once that scalar
condition is known, the functional part of the proof is independent of
the offspring law.

\begin{thebibliography}{SWX25}

\bibitem[A14]{A14}
E.~A\"id\'ekon.
\newblock Speed of the biased random walk on a Galton--Watson tree.
\newblock \emph{Probability Theory and Related Fields} \textbf{159}
(2014), 597--617.
\newblock \href{https://doi.org/10.1007/s00440-013-0515-y}{doi:10.1007/s00440-013-0515-y}.
\newblock Preprint: \href{https://arxiv.org/abs/1111.4313}{arXiv:1111.4313}.

\bibitem[BT20]{BT20}
A.~Bowditch and Y.~Tokushige.
\newblock Differentiability of the speed of biased random walks on
Galton--Watson trees.
\newblock \emph{ALEA, Latin American Journal of Probability and Mathematical
Statistics} \textbf{17} (2020), 609--642.
\newblock \href{https://doi.org/10.30757/ALEA.v17-24}{doi:10.30757/ALEA.v17-24}.
\newblock Preprint: \href{https://arxiv.org/abs/1906.07913}{arXiv:1906.07913}.

\bibitem[LPP96]{LPP96}
R.~Lyons, R.~Pemantle and Y.~Peres.
\newblock Biased random walks on Galton--Watson trees.
\newblock \emph{Probability Theory and Related Fields} \textbf{106}
(1996), 249--264.
\newblock \href{https://doi.org/10.1007/s004400050064}{doi:10.1007/s004400050064}.

\bibitem[MK26]{MK26}
M.~Mandarapu and S.~Kunkunuru.
\newblock A computer-assisted proof of speed monotonicity for the biased
random walk on a Galton--Watson tree beyond the known range.
\newblock Preprint (2026),
\href{https://arxiv.org/abs/2609.29894}{arXiv:2609.29894v1}.

\bibitem[SWX25]{SWX25}
H.~Song, L.~Wang and K.~Xiang.
\newblock The speed of a biased walk on a Galton--Watson tree without leaves
is monotonic for low values of bias.
\newblock \emph{Journal of Applied Probability} \textbf{62} (2025),
no.~3, 1044--1052.
\newblock \href{https://doi.org/10.1017/jpr.2024.113}{doi:10.1017/jpr.2024.113}.
\newblock Earlier version, \emph{Monotonicity of the speed for biased random
walk on Galton--Watson tree}:
\href{https://arxiv.org/abs/1610.08151}{arXiv:1610.08151v1}.
Theorem~1.2 and Lemma~2.5 cited in the text refer to this preprint's
numbering.

\end{thebibliography}
\end{document}
